% ============================================================
% PARTIE II - CHAPITRE II : Analyse et Conception
% ============================================================
\chapter{Analyse et Conception}
\pagestyle{styleGlobal}
\label{chap:analyse_conception}

% ============================================================
\section{Analyse du besoin}
% ============================================================

\subsection{Besoins Fonctionnels}

Les besoins fonctionnels représentent l'ensemble des fonctionnalités que «~SALISA~» doit offrir à ses utilisateurs. Nous les avons identifiés à partir de l'analyse des problèmes rencontrés dans la mise en relation manuelle entre prestataires et clients. Ces besoins sont les suivants :

\begin{itemize}
	\item inscription via un numéro de téléphone, avec validation par code OTP ;
	\item connexion et déconnexion sécurisées ;
	\item choix du rôle lors de l'inscription (prestataire, demandeur ou les deux) ;
	\item création et gestion du profil prestataire (informations, portfolio, tarifs, disponibilité, catégories de services) ;
	\item vérification d'identité et obtention de badges de certification ;
	\item publication de missions par le demandeur ;
	\item recherche et filtrage de prestataires par catégorie, arrondissement, disponibilité et budget ;
	\item mise en correspondance intelligente (matching) entre missions et prestataires ;
	\item messagerie directe entre demandeur et prestataire ;
	\item envoi et acceptation de devis structurés ;
	\item paiement sécurisé via Mobile Money (MTN MoMo et Airtel Money) \cite{mtn_momo} ;
	\item sécurisation des fonds par escrow jusqu'à la validation de la mission ;
	\item validation de la mission terminée et libération des fonds ;
	\item gestion des litiges entre prestataires et demandeurs ;
	\item notation et dépôt d'avis après chaque mission réalisée ;
	\item notifications multi-canaux (push PWA, WhatsApp, SMS) ;
	\item tableau de bord de modération pour les administrateurs.
\end{itemize}

\subsection{Besoins Techniques}

Les besoins techniques désignent l'ensemble des contraintes non fonctionnelles auxquelles un système doit se conformer, qu'il s'agisse de performances, de sécurité, de disponibilité ou de compatibilité avec l'environnement d'exécution. Pour «~SALISA~», ces contraintes sont directement liées au contexte congolais, et se déclinent de la façon suivante :

\begin{itemize}
	\item architecture Progressive Web Application (PWA) \cite{next_js}, mobile-first, utilisable depuis tous les terminaux sans installation ;
	\item compatibilité avec les réseaux à faible débit (2G/3G) pour garantir l'accès dans les zones peu couvertes ;
	\item fonctionnement hors ligne partiel pour la consultation des profils et des missions déjà chargés ;
	\item sécurité des échanges par protocoles HTTPS et WSS, authentification par JWT et OTP ;
	\item performance : temps de réponse des API inférieur à 2 secondes pour 95\% des requêtes ;
	\item disponibilité cible de 99,5\% ;
	\item intégration des API de paiement Mobile Money (MTN MoMo et Airtel Money) \cite{mtn_momo} ;
	\item base de données relationnelle (PostgreSQL) \cite{postgresql} avec cache Redis pour les sessions et les résultats fréquents ;
	\item messagerie en temps réel via WebSocket ;
	\item stockage des fichiers (photos de profil, pièces jointes) sur un service cloud externe.
\end{itemize}

% ============================================================
\section{Conception du système}
% ============================================================

\subsection{Spécifications fonctionnelles}

\subsubsection{Identification des acteurs}

Un acteur représente toute entité externe qui interagit avec le système «~SALISA~» \cite{jacobson_up}. Notre système comprend deux catégories d'acteurs : les acteurs humains et les acteurs systèmes.

\vspace{0.3cm}
\noindent\textbf{a- Acteurs humains}

\vspace{0.2cm}
Nous avons identifié quatre principaux acteurs humains :

\begin{itemize}
	\item \textbf{Visiteur :} utilisateur non encore enregistré sur la plateforme. Son unique interaction avec «~SALISA~» est l'inscription, qui lui permet de créer un compte et de devenir un utilisateur à part entière ;
	\item \textbf{Demandeur :} utilisateur connecté qui cherche et embauche des prestataires. Il publie des missions, contacte des prestataires, paie et note les prestations ;
	\item \textbf{Prestataire :} utilisateur connecté qui propose ses services. Il crée son profil, répond aux missions et envoie des devis ;
	\item \textbf{Administrateur :} utilisateur connecté chargé d'administrer la plateforme, de vérifier les identités, de gérer les litiges et de configurer le système.
\end{itemize}

\vspace{0.3cm}
\noindent\textbf{b- Acteurs systèmes}

\vspace{0.2cm}
En modélisation UML, on désigne par acteur système tout composant logiciel extérieur au système étudié qui interagit avec lui par l'échange de données ou de services. Nous en avons identifié deux pour «~SALISA~» :

\begin{itemize}
	\item \textbf{Service de paiement :} regroupe les APIs de MTN MoMo et d'Airtel Money. «~SALISA~» leur envoie des requêtes de paiement et en reçoit des confirmations de transaction \cite{mtn_momo} ;
	\item \textbf{Service de notifications :} regroupe Twilio SMS et WhatsApp Business API. «~SALISA~» lui envoie des instructions pour délivrer les codes OTP, les alertes et les messages aux utilisateurs.
\end{itemize}

\subsubsection{Diagramme de contexte statique}

Le diagramme de contexte statique présente «~SALISA~» au centre et montre toutes les entités externes qui interagissent avec lui. La figure~\ref{fig:contexte_statique} illustre ce diagramme.

\begin{figure}[H]
	\centering
	\includegraphics[width=0.95\textwidth]{diagramme_contexte_statique.png}
	\caption{Diagramme de contexte statique de «~SALISA~»}
	\label{fig:contexte_statique}
\end{figure}

Au centre se trouve le système «~SALISA~» représenté par un rectangle. Autour, les quatre acteurs humains (symbole stickman UML) interagissent avec lui : le Visiteur s'inscrit sur la plateforme ; le Demandeur publie des missions, paie et note ; le Prestataire crée son profil, répond et envoie des devis ; l'Administrateur administre et configure. Les deux acteurs systèmes (rectangles avec stéréotype \textit{$\langle\langle$système$\rangle\rangle$}) sont : le service de paiement pour les transactions, et le service de notifications pour les alertes SMS et WhatsApp.

\subsubsection{Identification des cas d'utilisation}

Un cas d'utilisation décrit une interaction entre un acteur et le système pour accomplir un objectif précis. Pour «~SALISA~», nous avons recensé 39 cas d'utilisation, regroupés en neuf familles selon les acteurs concernés. Les tableaux~\ref{tab:cu_f1} à~\ref{tab:cu_f9} présentent le détail de chacune de ces familles.

\begin{table}[H]
	\centering
	\begin{tabular}{|p{3cm}|p{10.5cm}|}
		\hline
		\multicolumn{2}{|c|}{\cellcolor{vertCongo}\textcolor{white}{\textbf{Famille 1 : Authentification et gestion du compte}}} \\
		\hline
		\rowcolor{vertCongo}
		\textcolor{white}{\textbf{Cas d'utilisation}} & \textcolor{white}{\textbf{Définition}} \\
		\hline
		\textbf{UC-01} & s'inscrire via numéro de téléphone (Visiteur) \\
		\hline
		\textbf{UC-02} & valider le code OTP (Visiteur) \\
		\hline
		\textbf{UC-03} & choisir son rôle lors de l'inscription (Visiteur) \\
		\hline
		\textbf{UC-04} & se connecter (Demandeur, Prestataire, Administrateur) \\
		\hline
		\textbf{UC-05} & se déconnecter (Demandeur, Prestataire, Administrateur) \\
		\hline
	\end{tabular}
	\caption{Cas d'utilisation de la famille F1 : Authentification et gestion du compte}
	\label{tab:cu_f1}
\end{table}

\begin{table}[H]
	\centering
	\begin{tabular}{|p{3cm}|p{10.5cm}|}
		\hline
		\multicolumn{2}{|c|}{\cellcolor{vertCongo}\textcolor{white}{\textbf{Famille 2 : Gestion du profil prestataire}}} \\
		\hline
		\rowcolor{vertCongo}
		\textcolor{white}{\textbf{Cas d'utilisation}} & \textcolor{white}{\textbf{Définition}} \\
		\hline
		\textbf{UC-06} & compléter son profil prestataire (Prestataire) \\
		\hline
		\textbf{UC-07} & gérer son portfolio (Prestataire) \\
		\hline
		\textbf{UC-08} & définir ses tarifs et sa disponibilité (Prestataire) \\
		\hline
		\textbf{UC-09} & gérer ses catégories et compétences (Prestataire) \\
		\hline
		\textbf{UC-10} & faire vérifier son identité pour obtenir le badge Vérifié (Prestataire, Administrateur) \\
		\hline
	\end{tabular}
	\caption{Cas d'utilisation de la famille F2 : Gestion du profil prestataire}
	\label{tab:cu_f2}
\end{table}

\begin{table}[H]
	\centering
	\begin{tabular}{|p{3cm}|p{10.5cm}|}
		\hline
		\multicolumn{2}{|c|}{\cellcolor{vertCongo}\textcolor{white}{\textbf{Famille 3 : Gestion des missions}}} \\
		\hline
		\rowcolor{vertCongo}
		\textcolor{white}{\textbf{Cas d'utilisation}} & \textcolor{white}{\textbf{Définition}} \\
		\hline
		\textbf{UC-11} & publier une mission (Demandeur) \\
		\hline
		\textbf{UC-12} & utiliser l'assistant IA pour la description (Demandeur) \\
		\hline
		\textbf{UC-13} & gérer ses missions publiées (Demandeur) \\
		\hline
	\end{tabular}
	\caption{Cas d'utilisation de la famille F3 : Gestion des missions}
	\label{tab:cu_f3}
\end{table}

\begin{table}[H]
	\centering
	\begin{tabular}{|p{3cm}|p{10.5cm}|}
		\hline
		\multicolumn{2}{|c|}{\cellcolor{vertCongo}\textcolor{white}{\textbf{Famille 4 : Recherche et matching}}} \\
		\hline
		\rowcolor{vertCongo}
		\textcolor{white}{\textbf{Cas d'utilisation}} & \textcolor{white}{\textbf{Définition}} \\
		\hline
		\textbf{UC-14} & rechercher un prestataire (Demandeur) \\
		\hline
		\textbf{UC-15} & filtrer les résultats de recherche (Demandeur) \\
		\hline
		\textbf{UC-16} & consulter le score de matching (Demandeur) \\
		\hline
		\textbf{UC-17} & recevoir des suggestions proactives de prestataires (Demandeur) \\
		\hline
		\textbf{UC-18} & consulter le profil d'un prestataire (Demandeur) \\
		\hline
	\end{tabular}
	\caption{Cas d'utilisation de la famille F4 : Recherche et matching}
	\label{tab:cu_f4}
\end{table}

\begin{table}[H]
	\centering
	\begin{tabular}{|p{3cm}|p{10.5cm}|}
		\hline
		\multicolumn{2}{|c|}{\cellcolor{vertCongo}\textcolor{white}{\textbf{Famille 5 : Messagerie et devis}}} \\
		\hline
		\rowcolor{vertCongo}
		\textcolor{white}{\textbf{Cas d'utilisation}} & \textcolor{white}{\textbf{Définition}} \\
		\hline
		\textbf{UC-19} & envoyer un message (Demandeur, Prestataire) \\
		\hline
		\textbf{UC-20} & envoyer un devis (Prestataire) \\
		\hline
		\textbf{UC-21} & accepter un devis (Demandeur) \\
		\hline
		\textbf{UC-22} & refuser ou demander une modification de devis (Demandeur) \\
		\hline
	\end{tabular}
	\caption{Cas d'utilisation de la famille F5 : Messagerie et devis}
	\label{tab:cu_f5}
\end{table}

\begin{table}[H]
	\centering
	\begin{tabular}{|p{3cm}|p{10.5cm}|}
		\hline
		\multicolumn{2}{|c|}{\cellcolor{vertCongo}\textcolor{white}{\textbf{Famille 6 : Paiement et escrow}}} \\
		\hline
		\rowcolor{vertCongo}
		\textcolor{white}{\textbf{Cas d'utilisation}} & \textcolor{white}{\textbf{Définition}} \\
		\hline
		\textbf{UC-23} & payer une mission via Mobile Money (Demandeur, service de paiement) \\
		\hline
		\textbf{UC-24} & choisir MTN MoMo comme moyen de paiement (Demandeur) \\
		\hline
		\textbf{UC-25} & choisir Airtel Money comme moyen de paiement (Demandeur) \\
		\hline
		\textbf{UC-26} & valider la mission terminée (Demandeur) \\
		\hline
		\textbf{UC-27} & scanner le code QR de validation sur place (Prestataire) \\
		\hline
		\textbf{UC-28} & ouvrir un litige (Demandeur, Prestataire) \\
		\hline
		\textbf{UC-29} & demander un retrait Mobile Money (Prestataire) \\
		\hline
	\end{tabular}
	\caption{Cas d'utilisation de la famille F6 : Paiement et escrow}
	\label{tab:cu_f6}
\end{table}

\begin{table}[H]
	\centering
	\begin{tabular}{|p{3cm}|p{10.5cm}|}
		\hline
		\multicolumn{2}{|c|}{\cellcolor{vertCongo}\textcolor{white}{\textbf{Famille 7 : Notation et réputation}}} \\
		\hline
		\rowcolor{vertCongo}
		\textcolor{white}{\textbf{Cas d'utilisation}} & \textcolor{white}{\textbf{Définition}} \\
		\hline
		\textbf{UC-30} & laisser un avis après mission (Demandeur, Prestataire) \\
		\hline
		\textbf{UC-31} & signaler un avis abusif (Demandeur, Prestataire) \\
		\hline
	\end{tabular}
	\caption{Cas d'utilisation de la famille F7 : Notation et réputation}
	\label{tab:cu_f7}
\end{table}

\begin{table}[H]
	\centering
	\begin{tabular}{|p{3cm}|p{10.5cm}|}
		\hline
		\multicolumn{2}{|c|}{\cellcolor{vertCongo}\textcolor{white}{\textbf{Famille 8 : Notifications}}} \\
		\hline
		\rowcolor{vertCongo}
		\textcolor{white}{\textbf{Cas d'utilisation}} & \textcolor{white}{\textbf{Définition}} \\
		\hline
		\textbf{UC-32} & recevoir des notifications multi-canaux (Demandeur, Prestataire, service de notifications) \\
		\hline
		\textbf{UC-33} & configurer ses préférences de notification (Demandeur, Prestataire) \\
		\hline
	\end{tabular}
	\caption{Cas d'utilisation de la famille F8 : Notifications}
	\label{tab:cu_f8}
\end{table}

\begin{table}[H]
	\centering
	\begin{tabular}{|p{3cm}|p{10.5cm}|}
		\hline
		\multicolumn{2}{|c|}{\cellcolor{vertCongo}\textcolor{white}{\textbf{Famille 9 : Administration}}} \\
		\hline
		\rowcolor{vertCongo}
		\textcolor{white}{\textbf{Cas d'utilisation}} & \textcolor{white}{\textbf{Définition}} \\
		\hline
		\textbf{UC-34} & accéder au tableau de bord de modération (Administrateur) \\
		\hline
		\textbf{UC-35} & valider ou rejeter une vérification d'identité (Administrateur) \\
		\hline
		\textbf{UC-36} & gérer les signalements (Administrateur) \\
		\hline
		\textbf{UC-37} & suspendre un compte utilisateur (Administrateur) \\
		\hline
		\textbf{UC-38} & configurer le système (Administrateur) \\
		\hline
		\textbf{UC-39} & consulter les statistiques avancées (Administrateur) \\
		\hline
	\end{tabular}
	\caption{Cas d'utilisation de la famille F9 : Administration}
	\label{tab:cu_f9}
\end{table}

\subsubsection{Relation entre les cas d'utilisation}

Les relations entre cas d'utilisation permettent de structurer et de clarifier les dépendances entre les fonctionnalités du système \cite{roques_uml}. On en distingue deux types principaux. La relation \textbf{$\langle\langle$include$\rangle\rangle$} indique qu'un cas d'utilisation est toujours exécuté dans le cadre d'un autre : par exemple, UC-01 (s'inscrire) inclut systématiquement UC-02 (valider le code OTP), car on ne peut pas créer de compte sans valider le code reçu par SMS ; de même, UC-21 (accepter un devis) inclut UC-23 (payer via Mobile Money), car l'acceptation d'un devis déclenche toujours le processus de paiement. La relation \textbf{$\langle\langle$extend$\rangle\rangle$}, quant à elle, indique qu'un cas d'utilisation est exécuté de façon optionnelle dans certaines conditions : par exemple, UC-11 (publier une mission) peut optionnellement étendre vers UC-12 (utiliser l'assistant IA), selon le choix du demandeur, et UC-26 (valider la mission) peut étendre vers UC-27 (scanner le code QR), uniquement pour les missions réalisées en présentiel.

\vspace{0.4cm}
La figure~\ref{fig:cu_global} présente le diagramme de cas d'utilisation global de «~SALISA~», organisé autour des neuf familles fonctionnelles.

\begin{figure}[H]
	\centering
	\includegraphics[width=0.95\textwidth]{diagramme_cu_global.png}
	\caption{Diagramme de cas d'utilisation global de «~SALISA~»}
	\label{fig:cu_global}
\end{figure}

Pour une meilleure lisibilité, chaque famille est ensuite détaillée dans un diagramme dédié, présentant les cas d'utilisation et leurs relations avec précision. Les figures~\ref{fig:cu_f1} à~\ref{fig:cu_f9} illustrent ces neuf diagrammes.

\begin{figure}[H]
	\centering
	\includegraphics[width=0.85\textwidth]{diagramme_cu_f1.png}
	\caption{Diagramme de cas d'utilisation - F1 : Authentification}
	\label{fig:cu_f1}
\end{figure}

\begin{figure}[H]
	\centering
	\includegraphics[width=0.85\textwidth]{diagramme_cu_f2.png}
	\caption{Diagramme de cas d'utilisation - F2 : Profil prestataire}
	\label{fig:cu_f2}
\end{figure}

\begin{figure}[H]
	\centering
	\includegraphics[width=0.85\textwidth]{diagramme_cu_f3.png}
	\caption{Diagramme de cas d'utilisation - F3 : Gestion des missions}
	\label{fig:cu_f3}
\end{figure}

\begin{figure}[H]
	\centering
	\includegraphics[width=0.85\textwidth]{diagramme_cu_f4.png}
	\caption{Diagramme de cas d'utilisation - F4 : Recherche et matching}
	\label{fig:cu_f4}
\end{figure}

\begin{figure}[H]
	\centering
	\includegraphics[width=0.85\textwidth]{diagramme_cu_f5.png}
	\caption{Diagramme de cas d'utilisation - F5 : Messagerie et devis}
	\label{fig:cu_f5}
\end{figure}

\begin{figure}[H]
	\centering
	\includegraphics[width=0.85\textwidth]{diagramme_cu_f6.png}
	\caption{Diagramme de cas d'utilisation - F6 : Paiement et escrow}
	\label{fig:cu_f6}
\end{figure}

\begin{figure}[H]
	\centering
	\includegraphics[width=0.85\textwidth]{diagramme_cu_f7.png}
	\caption{Diagramme de cas d'utilisation - F7 : Notation et réputation}
	\label{fig:cu_f7}
\end{figure}

\begin{figure}[H]
	\centering
	\includegraphics[width=0.85\textwidth]{diagramme_cu_f8.png}
	\caption{Diagramme de cas d'utilisation - F8 : Notifications}
	\label{fig:cu_f8}
\end{figure}

\begin{figure}[H]
	\centering
	\includegraphics[width=0.85\textwidth]{diagramme_cu_f9.png}
	\caption{Diagramme de cas d'utilisation - F9 : Administration}
	\label{fig:cu_f9}
\end{figure}

\subsubsection{Classification des cas d'utilisation (priorités, risques et itérations)}

Le tableau~\ref{tab:classification_cu} présente la classification des principaux cas d'utilisation selon leur priorité fonctionnelle, leur niveau de risque technique et l'itération de développement dans laquelle ils s'inscrivent. L'itération 1 regroupe les fonctionnalités essentielles au fonctionnement minimal du système. L'itération 2 contient les fonctionnalités importantes mais non bloquantes. L'itération 3 rassemble les fonctionnalités souhaitables à plus long terme.

\begin{table}[H]
	\centering
	\begin{tabular}{|c|p{4.5cm}|c|c|c|}
		\hline
		\rowcolor{vertCongo}
		\textcolor{white}{\textbf{UC}} &
		\textcolor{white}{\textbf{Libellé}} &
		\textcolor{white}{\textbf{Priorité}} &
		\textcolor{white}{\textbf{Risque}} &
		\textcolor{white}{\textbf{Itération}} \\
		\hline
		UC-01 & s'inscrire & Élevée & Élevé & 1 \\
		\hline
		UC-02 & valider le code OTP & Élevée & Élevé & 1 \\
		\hline
		UC-03 & choisir son rôle & Élevée & Faible & 1 \\
		\hline
		UC-06 & compléter son profil & Élevée & Moyen & 1 \\
		\hline
		UC-11 & publier une mission & Élevée & Moyen & 1 \\
		\hline
		UC-14 & rechercher un prestataire & Élevée & Moyen & 1 \\
		\hline
		UC-19 & envoyer un message & Élevée & Moyen & 1 \\
		\hline
		UC-20 & envoyer un devis & Élevée & Moyen & 1 \\
		\hline
		UC-23 & payer via Mobile Money & Élevée & Élevé & 1 \\
		\hline
		UC-26 & valider la mission & Élevée & Élevé & 1 \\
		\hline
		UC-30 & laisser un avis & Élevée & Faible & 1 \\
		\hline
		UC-32 & recevoir des notifications & Élevée & Moyen & 1 \\
		\hline
		UC-34 & tableau de bord admin & Élevée & Moyen & 1 \\
		\hline
		UC-12 & utiliser l'assistant IA & Moyenne & Élevé & 2 \\
		\hline
		UC-16 & score de matching & Moyenne & Élevé & 2 \\
		\hline
		UC-28 & ouvrir un litige & Moyenne & Moyen & 2 \\
		\hline
		UC-38 & configurer le système & Basse & Faible & 3 \\
		\hline
	\end{tabular}
	\caption{Classification des principaux cas d'utilisation de «~SALISA~»}
	\label{tab:classification_cu}
\end{table}

\subsection{Spécification détaillée}

\subsubsection{Description textuelle des cas d'utilisation (au plus 2 cu)}

\noindent\textbf{a- Description de UC-01 : s'inscrire via numéro de téléphone}

\vspace{0.2cm}
\begin{table}[H]
	\centering
	\begin{tabular}{|p{3.5cm}|p{10cm}|}
		\hline
		\rowcolor{vertClair}
		\textbf{Cas d'utilisation} & UC-01 : s'inscrire via numéro de téléphone \\
		\hline
		\textbf{Acteur principal} & Visiteur \\
		\hline
		\textbf{Précondition} & le visiteur n'a pas encore de compte sur «~SALISA~» \\
		\hline
		\textbf{Déclencheur} & le visiteur clique sur le bouton «~S'inscrire~» \\
		\hline
		\textbf{Scénario nominal} &
		\begin{enumerate}
			\item Le visiteur saisit son numéro de téléphone au format +242 XXXXXXXXX.
			\item Le système valide le format et envoie un code OTP par SMS (UC-02).
			\item Le visiteur saisit le code OTP reçu.
			\item Le système vérifie le code (durée de validité : 5 minutes).
			\item Le visiteur choisit son rôle : prestataire, demandeur, ou les deux (UC-03).
			\item Le compte est créé et le visiteur devient un utilisateur connecté.
		\end{enumerate} \\
		\hline
		\textbf{Scénarios alternatifs} &
		\begin{itemize}
			\item Code OTP expiré : le visiteur peut en demander un nouveau.
			\item 3 tentatives échouées : blocage temporaire de 10 minutes.
			\item Numéro déjà existant : redirection vers la page de connexion.
		\end{itemize} \\
		\hline
		\textbf{Postcondition} & le compte est créé ; le visiteur est devenu un utilisateur et peut accéder à toutes les fonctionnalités selon son rôle \\
		\hline
	\end{tabular}
	\caption{Description textuelle de UC-01 : s'inscrire via numéro de téléphone}
	\label{tab:uc01}
\end{table}

\vspace{0.5cm}
\noindent\textbf{b- Description de UC-11 : publier une mission}

\vspace{0.2cm}
\begin{table}[H]
	\centering
	\begin{tabular}{|p{3.5cm}|p{10cm}|}
		\hline
		\rowcolor{vertClair}
		\textbf{Cas d'utilisation} & UC-11 : publier une mission \\
		\hline
		\textbf{Acteur principal} & Demandeur \\
		\hline
		\textbf{Précondition} & l'utilisateur est connecté en tant que demandeur \\
		\hline
		\textbf{Déclencheur} & le demandeur clique sur «~Publier une mission~» \\
		\hline
		\textbf{Scénario nominal} &
		\begin{enumerate}
			\item Le demandeur saisit le titre de la mission.
			\item Il renseigne la description de son besoin (ou utilise l'assistant IA - UC-12).
			\item Il sélectionne la catégorie et l'arrondissement.
			\item Il indique le budget et le délai souhaité.
			\item Il confirme et publie la mission.
			\item Le système lance l'algorithme de matching et notifie les prestataires correspondants.
		\end{enumerate} \\
		\hline
		\textbf{Scénarios alternatifs} &
		\begin{itemize}
			\item Le demandeur n'est pas connecté : redirection vers la page de connexion.
			\item Aucun prestataire ne correspond : message «~Aucun résultat proche~».
		\end{itemize} \\
		\hline
		\textbf{Postcondition} & la mission est publiée et les prestataires correspondants sont notifiés \\
		\hline
	\end{tabular}
	\caption{Description textuelle de UC-11 : publier une mission}
	\label{tab:uc11}
\end{table}

\subsubsection{Diagramme de séquence (au plus 2)}

La figure~\ref{fig:seq_inscription} illustre le diagramme de séquence du cas d'utilisation UC-01 (s'inscrire via numéro de téléphone).

\begin{figure}[H]
	\centering
	\includegraphics[width=0.95\textwidth]{diagramme_sequence_inscription.png}
	\caption{Diagramme de séquence - inscription via numéro de téléphone (UC-01)}
	\label{fig:seq_inscription}
\end{figure}

La figure~\ref{fig:seq_paiement} illustre le diagramme de séquence du scénario de paiement et de libération des fonds (UC-23 et UC-26).

\begin{figure}[H]
	\centering
	\includegraphics[width=0.95\textwidth]{diagramme_sequence_paiement.png}
	\caption{Diagramme de séquence - paiement et libération des fonds (UC-23 et UC-26)}
	\label{fig:seq_paiement}
\end{figure}

\subsubsection{Diagramme d'activités (au plus 2)}

La figure~\ref{fig:act_publication} présente le diagramme d'activités du processus de publication d'une mission.

\begin{figure}[H]
	\centering
	\includegraphics[width=0.95\textwidth]{diagramme_activites_mission.png}
	\caption{Diagramme d'activités - publication de mission et sélection du prestataire}
	\label{fig:act_publication}
\end{figure}

La figure~\ref{fig:act_paiement} illustre le diagramme d'activités du processus de paiement sécurisé par escrow.

\begin{figure}[H]
	\centering
	\includegraphics[width=0.95\textwidth]{diagramme_activites_paiement.png}
	\caption{Diagramme d'activités - paiement sécurisé par escrow}
	\label{fig:act_paiement}
\end{figure}

\subsection{Réalisation des cas d'utilisation}

\subsubsection{Quelques règles de gestion}

Les règles de gestion définissent le comportement du système «~SALISA~» dans des situations précises \cite{roques_uml}. Elles garantissent la cohérence des données et la fiabilité des transactions.

\vspace{0.2cm}
\noindent\textbf{Gestion des comptes}

\begin{itemize}
	\item RG-001 : un numéro de téléphone ne peut être associé qu'à un seul compte ;
	\item RG-002 : un même compte peut avoir simultanément les rôles de prestataire et de demandeur ;
	\item RG-003 : un compte est suspendu automatiquement après trois signalements validés par un administrateur ;
	\item RG-004 : lors de la suppression d'un compte, les données sont conservées 30 jours avant effacement définitif.
\end{itemize}

\noindent\textbf{Gestion des missions et des prestataires}

\begin{itemize}
	\item RG-010 : un profil prestataire n'est visible dans la recherche qu'à partir de 70\% de complétion ;
	\item RG-011 : un prestataire avec un compte gratuit ne peut postuler qu'à 3 missions par mois ;
	\item RG-012 : une mission sans activité depuis 30 jours est archivée automatiquement ;
	\item RG-013 : au maximum 10 prestataires sont notifiés par mission (les 10 meilleurs scores de matching).
\end{itemize}

\noindent\textbf{Gestion des paiements}

\begin{itemize}
	\item RG-020 : la commission de «~SALISA~» est de 8\% du montant total de chaque transaction ;
	\item RG-021 : le montant minimum d'une transaction est de 2 000 FCFA ;
	\item RG-022 : si le demandeur ne valide pas la mission dans les 7 jours suivant la livraison, les fonds sont libérés automatiquement au prestataire ;
	\item RG-023 : un litige doit être ouvert dans les 48 heures suivant la date de livraison prévue.
\end{itemize}

\noindent\textbf{Gestion des avis}

\begin{itemize}
	\item RG-030 : un avis ne peut être laissé qu'après une mission avec transaction validée ;
	\item RG-031 : le délai pour laisser un avis est de 14 jours après la validation de la mission ;
	\item RG-032 : les avis sont définitifs et ne peuvent jamais être supprimés par le prestataire ou le demandeur.
\end{itemize}

\subsubsection{Modèle du domaine (Diagramme de classe du système)}

La figure~\ref{fig:diag_classes} présente le diagramme de classes de «~SALISA~», qui représente les entités principales du système, leurs attributs, leurs méthodes et les relations qui les unissent \cite{booch_uml}.

\begin{figure}[H]
	\centering
	\includegraphics[width=0.95\textwidth]{diagramme_classes.png}
	\caption{Diagramme de classes du système «~SALISA~»}
	\label{fig:diag_classes}
\end{figure}

\subsubsection{Diagramme d'Objets}

La figure~\ref{fig:diag_objets} illustre un exemple d'instanciation du système «~SALISA~», correspondant au scénario où un demandeur publie une mission, un prestataire y répond avec un devis, et le paiement est sécurisé par escrow.

\begin{figure}[H]
	\centering
	\includegraphics[width=0.95\textwidth]{diagramme_objets.png}
	\caption{Diagramme d'objets - scénario de publication et réponse à une mission}
	\label{fig:diag_objets}
\end{figure}

\subsection{Conception architecturale}

\subsubsection{Architecture logicielle}

«~SALISA~» repose sur une architecture modulaire organisée autour de cinq couches principales.

Le \textbf{front-end} est développé en Next.js 14 \cite{next_js}, un framework React qui permet de générer des pages côté serveur (SSR) et des pages statiques (SSG). Il assure l'aspect PWA mobile-first et l'optimisation pour les réseaux lents.

Le \textbf{back-end} est développé en Node.js avec Express.js. Il expose une API REST pour toutes les opérations classiques et utilise Socket.io pour la messagerie en temps réel via WebSocket.

La \textbf{base de données} principale est PostgreSQL \cite{postgresql}, qui stocke toutes les données persistantes (utilisateurs, missions, transactions, avis). Redis est utilisé comme cache pour les sessions et les résultats de recherche fréquents.

Le \textbf{service de matching} est un microservice Python développé avec FastAPI. Il calcule les scores de compatibilité entre les missions et les prestataires sur la base de huit critères pondérés : localisation, disponibilité, compétences, budget, note, fiabilité, temps de réponse et niveau de badge.

Le \textbf{service de notifications} est un worker asynchrone qui envoie les alertes via trois canaux : les notifications push PWA, les messages WhatsApp et les SMS.

\subsubsection{Architecture globale de la solution (Diagramme de déploiement)}

La figure~\ref{fig:diag_deploiement} présente le diagramme de déploiement de «~SALISA~», qui montre la répartition des composants sur les équipements matériels et les services en ligne.

\begin{figure}[H]
	\centering
	\includegraphics[width=0.95\textwidth]{diagramme_deploiement.png}
	\caption{Diagramme de déploiement de «~SALISA~»}
	\label{fig:diag_deploiement}
\end{figure}